
\subsubsection{6.1 Key Findings}

\textbf{Benchmark saturation is a phase transition, not a gradual decline.} Prior work modeled saturation as a smooth monotonic trend (Liao et al. 2022, S\_index); the results here demonstrate it is a discrete structural break detectable by standard change-point analysis. A threshold supports decision-making (flag benchmarks above it for review); a trend does not.

\textbf{The two regimes are quantitatively extreme.} Pre-breakpoint variance (3.$81\times$ global) and post-breakpoint variance (0.$20\times$ pre-segment) differ by approximately $19\times$ (3.814 / 0.1981 $\approx$ 19.25). This is visible in the raw scatter (Figure 2) --- not a subtle statistical artifact.

\textbf{Three experiments with distinct test designs corroborate the same mechanism.} H-E1 detects the break; H-M1 confirms the pre-breakpoint character; H-M2 confirms the post-breakpoint character. Each uses pre-specified metrics and independent statistical tests, though all three analyze the same underlying dataset and breakpoint.

\textbf{Interpretation of OLS reversal (rho: $-0.28 \rightarrow$ +0.137).} The reversal is not a contradiction --- but its cause is not established with certainty. The most plausible interpretation is PwC benchmark database composition changes between snapshots (N=111 $\rightarrow$ N=115): if newly added benchmarks cluster at high paper\_count with high CoV, they would locally reverse the aggregate trend slope. This cannot be confirmed without per-snapshot benchmark identity data. PELT's residual-based approach is insulated from this reversal regardless of its origin.

\subsubsection{6.2 Limitations}

\textbf{L1: Small pre-segment (n\_pre=8) limits directional characterization.} The structural break falls near the left boundary of the paper\_count distribution. Moment-based inference (skewness, H-M3 M1/M3) is unreliable at n=8. Primary claims (structural break, variance compression) rely on statistics robust to small segment sizes. This is a consequence of where the natural breakpoint falls, not a methodological flaw.

\textbf{L2: Cross-sectional design cannot establish causality.} Benchmarks are pooled at a single time point. The two-regime structure is consistent with Goodhart saturation dynamics, but cross-sectional analysis cannot rule out alternative explanations (benchmark selection effects, metric heterogeneity across task types, benchmark age confounds). Longitudinal within-benchmark CoV trajectories would be required for causal attribution.

\textbf{L3: Dataset snapshot sensitivity.} paper\_count\*=39 reflects the Aug 2026 PwC snapshot. The OLS trend direction already reversed between Phase 1 (N=111) and the current analysis (N=115), demonstrating snapshot sensitivity. Periodic re-analysis is recommended.

\textbf{L4: Bootstrap CI width = 31.5.} The 95\% CI [38, 69.5] exceeds the soft $\leq 20-paper$ criterion, driven by the small pre-segment. The point estimate (39) is actionable; the CI should be reported in any application.

\textbf{L5: Domain stratification not tested.} A global paper\_count\* is reported. Whether task types (image classification, NLP) have different saturation thresholds is an explicit open question.

\textbf{L6: Multiple testing.} Three primary hypotheses (H-E1, H-M1, H-M2) are tested with raw p-values. Under Bonferroni correction ($\alpha /3 \approx$ 0.017), the H-E1 permutation p=0.035 does not survive correction; the corroborating piecewise F-test (p=0.0021) does. The permutation test and piecewise F-test are treated as two convergent lines of evidence rather than independent gates. Uncorrected values are reported; the permutation p=0.035 should not be interpreted as a standalone significance claim.

\textbf{L7: PELT L2 cost function assumption.} The L2 cost assumes Gaussian residuals. Whether benchmark residual CoV follows a Gaussian distribution is not formally tested. The robustness sweep across penalties (Figure 4) confirms the single breakpoint is stable, but non-Gaussian distributions (e.g., heavy-tailed CoV) could affect PELT's detection sensitivity.

\subsubsection{6.3 Broader Impact}

This work provides a methodology for operationalizing benchmark retirement decisions using existing leaderboard data. The threshold paper\_count\*$\approx 39$ should be interpreted as a benchmark health indicator, not a hard retirement trigger. Benchmarks above this threshold merit closer review including rank reversal rates, standard deviation of top-k methods, and qualitative assessment of whether improvements represent genuine methodological advances.

Potential misuse: mechanically retiring benchmarks above paper\_count\*=39 without qualitative review could prematurely sunset benchmarks where post-saturation reflects genuine consensus rather than Goodhart gaming. Retirement decisions should combine quantitative indicators with community input.

